\begin{textsample}{POS Dim 6 – global\_north\_2025\_03 – Score 5.00 – global\_north\_2025\_03/122728904.txt}  \label{ex:f6_pos_228}
Students set up camp on campus and won't leave until demands are met This article is brought to you by our \textbf{exclusive} \textbf{subscriber} partnership with our sister title USA Today, and has been written by our American colleagues.
It does not necessarily reflect the view of The Herald.
Students have begun camping on University Avenue as part of their protest against Glasgow University ' s ties with the arms industry and Israel.
The demonstrators have pitched their tents, vowing not to leave until the university meets their demands.
Hundreds of students rallied outside the University of Glasgow library at 2pm today, March 24 in support of five students now on their sixth day of a hunger strike.
The protest stems from a campaign by the Glasgow University Justice for Palestine Society ( GUJPS ), which seeks the university ' s complete dissociation from arms companies and an academic boycott of Israel.
They claim the university ' s? 6.8 million investment in arms manufacturers, including BAE Systems, implicates it in war crimes, @ @ @ @ @ @ @ @ @ @ in a hunger strike in response to the universities ties ( Image : Gordon Terris ) A Palestinian student participating in the encampment said :"I ' m here today because Glasgow University can not keep profiting off the blood of my people in Palestine with immunity."As students of \textbf{conscience}, we ' re here to hold this university accountable and show them that no repressive action they take will deter us from fighting for Palestinian liberation."Another student highlighted the university ' s significant investment in arms, saying :"They invest over? 6 million in the murderous arms industry -- including in BAE Systems, who make missile systems for F-35 fighter jets, the same jets Israel used in its genocidal bombardment of Palestinians in Gaza."Enough is enough."We will not leave until the university makes a firm, public commitment to divest from these \textbf{merchants} of death."( Image : Gordon Terris ) University rector Dr Ghassan Abu-Sittah stood in solidarity with the students, stating :"I @ @ @ @ @ @ @ @ @ @ genocide and UK Government complicity."The slaughter of innocent civilians is being carried out using weapons that are partly or fully manufactured by UK arms companies that the university holds shares in and is thus a part owner of."Last Wednesday, March 19, around 25 students from GUJPS occupied the Charles Wilson Building to demand the university sever ties with the arms industry.
The university responded by calling the police to campus, who blocked the entrance to the building, prevented food and water from reaching those inside, and threatened to arrest the students involved.
The students are refusing to move until their demands are met ( Image : Gordon Terris ) Commenting on these events, one of the student activists said :"The police have been called onto campus every day for the past three days."This represents a dramatic escalation in how our university manages student protests and indicates their refusal to engage with their students and staff on the vital issue of divestment."Furthermore, @ @ @ @ @ @ @ @ @ @ their lack of care for the safety and welfare of students."Glasgow University student Hannah Taylor has been banned from campus following a non-violent protest with Youth Demand, where they sprayed water-soluble paint over a university building.
This banning order means that they are unable to attend their lectures and may fail their degree before their disciplinary hearing.
The university rector and trauma surgeon, Dr Ghassan Abu-Sittah, described the decision as"disproportionate and authoritarian"and called for the ban to be reversed.
Recent protests at the university led them to launch a consultation on their socially responsible investment policy ( SRI ) which was open to all students and staff.
The results of the consultation showed that 70 per cent of respondents believed that the university should always exclude investments in defence firms.
Despite these results, last semester the university court voted for a new SRI policy which did not include divestment from arms companies.
This decision was criticised by the Student Representative Council ( SRC ), who @ @ @ @ @ @ @ @ @ @ A University spokesperson said :"The University of Glasgow upholds the right to freedom of expression, including the right of staff and students to engage in peaceful demonstrations.
However, we do not tolerate activities which interfere with the rights of others to go about their business in peace."As an institution, we stand against hate or harassment of any kind.
We regularly communicate with all our staff and students about the need for tolerance towards each other, and we reiterate this call for all members of our community to be respectful to each other at all times."The University has been in contact with GUJPS on numerous occasions and has specifically offered assistance and support to those saying they are on hunger strike."Comments \&amp ; Moderation Readers ' comments : You are personally liable for the content of any comments you upload to this website, so please act responsibly.
We do not pre-moderate or monitor readers ' comments appearing on our websites, but we do post-moderate in response to complaints @ @ @ @ @ @ @ @ @ @ our attention.
You can make a complaint by using the ' report this post ' link.
We may then apply our discretion under the user terms to amend or delete comments.
Post moderation is undertaken full-time 9am-6pm on weekdays, and on a part-time basis outwith those hours. follow us This website and associated newspapers adhere to the Independent Press Standards Organisation ' s Editors ' Code of Practice.
If you have a complaint about the editorial content which relates to inaccuracy or intrusion, then please contact the editor here.
If you are \textbf{dissatisfied} with the response provided you can contact IPSO here

% matched lemmas: conscience, dissatisfy, exclusive, merchant, subscriber
\end{textsample}
